\scanpage{51}
\mspage{\handwritten{19}}
% archive sheet number: 50

\deleted{ouvert (EGA IV 2.4.9), il \deleted{\unsure{pas}} suffit de prouver qu'i\deleted{o}\inserted{l} existe une telle extension $k'$ et un point $x'$ de $H'=H\otimes_k k'$ au dessus de $x$, tels que $H'-G'$ soit un voisinage de $x'$. Or \deleted{on}\handwritten{il} suffira de prendre \deleted{\unsure{et}} $k'=k(x)$, de sorte que $G'$ admet un point $x'$ rationnel sur $k'$ au dessus de $x$. Alors l'ouvert dans $H'$ translaté de $G'$ par $x'$ contient $x'$ et est contenu dans $H'-G'$\handwritten{,}\deleted{\unsure{ce} qu} qui est donc un voisinage de $x'$, ce qui achève la démonstration de 11.11.}

\uline{Remarque} 11.11.\deleted{\unsure{3|2}}\handwritten{1}. \deleted{Monsieur} M. Raynaud nous signale qu'on peut générali-\linebreak ser 11.11 au cas d'un schéma en groupes quasi-compact et quasi-s\deleted{p}\handwritten{é}\deleted{é}\handwritten{p}aré $G$ sur un schéma\inserted{$S$ localement n\oe{} thérien} régulier de dimension $\handwritten{\leqslant}1$, en supposant de plus $G$ \uline{plat} sur $S$. O$\overset{\handwritten{-\,\cdot}}{\underset{\handwritten{n}}{\text{n}}}$ a alors équivalence des conditions suivantes:

(i) $G$ est affine sur $S$.

(ii) $G$ est quasi-affine sur $S$.

(iii) Gn peut faire opérer $G$ fidèlement sur un $S$-schéma quasi-affine.

(iv) On peut faire opérer $G$ fidèlement sur un Module quasi-cohérent plat sur $S$.

(v) (Si $G$ est de type fini sur $S$\inserted{et $S$ noethérien\handwritten{\deleted{\unsure{?}}.}}) $G$ est isomorphe à un sous-groupe fermé d'un $\text{\uline{Aut}}_V$, où $V$ est un Module localement libre de type fini sur $\deleted{\unsure{S|B}}\handwritten{S}$.

\margin{\figure{\handwritten{\begin{tikzpicture}[x=0.35cm,y=0.35cm]\draw (1,0) .. controls (0.1,-0.3) and (0.1,0.1) .. (0.1,0.7) -- (0.1,2) .. controls (0.1,2.7) and (0.4,2.8) .. (1,2.8);\draw (0.2,2.5) .. controls (-.6,3.5) and (-.7,6.2) .. (-.9,8.0);\draw[->] (-.92,8.15) .. controls (-.97,9.3) and (-1,10.5) .. (-.8,11) .. controls (-.4,12) and (0.5,12.4) .. (1.6,12.5);\end{tikzpicture}}}}
\deleted{\uline{Remarque} 11.11.2. Avec les notations de la démonstration précédente, on voit facilement (une fois qu'on sait que $G$ est affine) que \unsure{\deleted{G}l|\deleted{l}G}'homomor-\linebreak phisme $G\handwritten{\longrightarrow}G'=\underset{\handwritten{\longleftarrow}}{\lim}G_i$ est un isomorphisme, ce qui donne pour $G$ un théorème de structure\handwritten{,} qui sera \deleted{repris}\inserted{prouvé} par une autre méthode dans 11.1\deleted{\unsure{4|2}}\handwritten{\unsure{5|3}}.}
